\assignq{3}{4}{Land\'e factors: electronic and nuclear coupling}{2}
\subsection*{(a) Derive \(g_J\)}
Write \(\mathbf J=\mathbf L+\mathbf S\),
\(\boldsymbol\mu=-\mu_B(g_L\mathbf L+g_S\mathbf S)\)
in dimensionless angular-momentum units,
with \(g_L=1,\ g_S=2\).
Only the projection of \(\boldsymbol\mu\)
on \(\mathbf J\) contributes to the
weak-field diagonal energy in a fixed \(J\) manifold:
\[
g_J=\frac{\langle(g_L\mathbf L+g_S\mathbf S)
\cdot\mathbf J\rangle}{J(J+1)}.
\]
Using
\[
2\mathbf J\cdot\mathbf L
=\mathbf J^2+\mathbf L^2-\mathbf S^2,\quad
2\mathbf J\cdot\mathbf S
=\mathbf J^2+\mathbf S^2-\mathbf L^2,
\]
we find
\[
\boxed{g_J=1+
\frac{J(J+1)+S(S+1)-L(L+1)}{2J(J+1)}}
\quad (J>0).
\]

\subsection*{(b) Derive \(g_F\) neglecting nuclear magnetism}
Now \(\mathbf F=\mathbf I+\mathbf J\).
The electronic moment lies along \(\mathbf J\)
on average, while \(\mathbf J\) projects onto
\(\mathbf F\) with factor
\[
\frac{\langle\mathbf J\cdot\mathbf F\rangle}{F(F+1)}
=\frac{F(F+1)+J(J+1)-I(I+1)}{2F(F+1)}.
\]
Thus
\[
\boxed{g_F=g_J
\frac{F(F+1)+J(J+1)-I(I+1)}{2F(F+1)}},
\qquad F>0.
\]
This is the electronic contribution only;
the nuclear \(g_I\mu_N\) term was excluded
by the question's stated approximation.

\subsection*{(c) \({}^2P_{1/2}\) in a weak field}
Here \(L=1,\ S=1/2,\ J=1/2\) and
\[
g_J=1+\frac{3/4+3/4-2}{2(3/4)}
=\boxed{\frac23}.
\]
For \(I>0\), the coupled manifolds are
\(F_+=I+\tfrac12\) and \(F_-=I-\tfrac12\)
(the latter is \(F=0\) for \(I=1/2\)).
Substituting into the preceding expression
gives, for nonzero \(F\),
\[
\boxed{g_{I+1/2}=+\frac{2}{3(2I+1)},\qquad
g_{I-1/2}=-\frac{2}{3(2I+1)}}.
\]
The weak-field energies are
\[
\boxed{\Delta E_{Fm_F}=\mu_B g_Fm_FB},
\quad m_F=-F,\ldots,F.
\]
The upper/lower hyperfine multiplets
fan out into \(2F+1\) equally spaced
Zeeman levels with \emph{opposite} slopes.
For the common \(I=5/2\) example,
\(F=3,2\) with \(g_F=+1/9,-1/9\).
When \(F=0\), there is only \(m_F=0\)
and no first-order electronic Zeeman
splitting, so the formula defining \(g_{F=0}\)
should not be used.
\methodnote{Project angular momenta
sequentially \(L,S\to J\), then \(I,J\to F\).
In strong fields this coupled-\(F\) result
no longer applies without state mixing.}
